% SOURCE_AUTHORITY: sga5_fr_workpass.tex, lines 7229--7266 (LF-normalized unit).
% SOURCE_SHA256: 791F4EFFC5E02832D5D77ED03518C8156D6F07E4C8238B03545DB93D883FBB28
\subsection*{4. Filtraciones y graduaciones}
\subsubsection*{4.1. Generalidades}
\subsubsection*{4.1.1}
Sean $\cC$ una categoría abeliana y $X=(X_n)_{n\in\ZZ}$ un sistema proyectivo indexado por $\ZZ$ con valores en $\cC$. Se supone
\[
X_i=0\quad(i<0).
\]
Para todo elemento $p$ de $\ZZ$, se define un subobjeto $F^pX$ de $X$ en la categoría $P_1=\Hom(\ZZ^\circ,\cC)$ poniendo
\[
(F^pX)_n=F^pX_n=
\begin{cases}
\Ker(X_n\to X_{p-1}),& n\ge p-1,\\
0,& n<p-1,
\end{cases}
\]
con los morfismos de transición definidos de manera evidente. Se obtiene así una filtración decreciente $(F^pX)_{p\in\ZZ}$ de $X$ que verifica:
\begin{equation}\tag{a}
F^pX=X\quad\text{si }p\le0,
\end{equation}
\begin{equation}\tag{b}
(F^pX)_{p-1}=0\quad\text{para todo }p.
\end{equation}
Por otra parte, el lema de la serpiente aplicado al diagrama conmutativo exacto $(D)$ que aparece a continuación demuestra que, siempre que $n\ge p$, el morfismo canónico
\[
F^pX_n/F^{p+1}X_n\longrightarrow \Im(X_n\to X_p)\cap\Ker(X_p\to X_{p-1})
\]
es un isomorfismo:
\[
\begin{tikzcd}[column sep=large,row sep=large]
0 \arrow[r] & \Ker(X_n\to X_p) \arrow[r] \arrow[d] & X_n \arrow[r] \arrow[d] & \Im(X_n\to X_p) \arrow[r] \arrow[d] & 0\\
0 \arrow[r] & 0 \arrow[r] & X_{p-1} \arrow[r,"\sim"'] & X_{p-1} \arrow[r] & 0 .
\end{tikzcd}
\]
Si ahora $X$ es un sistema proyectivo indexado por $\NN$ con valores en $\cC$, se prolonga a un sistema proyectivo indexado por $\ZZ$ poniendo $X_n=0$ $(n<0)$, con los morfismos de transición evidentes, lo que permite aplicarle las construcciones precedentes. Esto nos permite considerar asimismo la categoría
\[
P=\Hom(\NN^\circ,\cC)
\]
como una subcategoría plena de $P_1$.
